\phead{SCHMIDT AND DENSITY MATRICES}{5a. Why orthogonality matters}{Schmidt form reveals the spectra of both reduced states.}
A Schmidt decomposition is a coherent sum of paired, orthonormal local vectors:
\[
\ket\Psi=\sum_k\sqrt{\lambda_k}\ket{a_k}\ket{b_k},
\qquad \lambda_k>0,\quad\sum_k\lambda_k=1.
\]
Its joint density matrix contains coherence between different branch labels:
\[
\rho_{AB}=\sum_{j,k}\sqrt{\lambda_j\lambda_k}\,
\ket{a_j}\bra{a_k}\otimes\ket{b_j}\bra{b_k}.
\]
The whole state is pure: the joint spectrum is \((1,0,\ldots)\). Taking a partial trace describes access to one subsystem; it does not measure the partner or destroy the joint cross terms.

\lead{Orthogonality has a separate role on each side}
For any paired expansion of this form, before assuming orthogonality,
\[
\rho_A=\sum_{j,k}\sqrt{\lambda_j\lambda_k}\,
\braket{b_k|b_j}\ket{a_j}\bra{a_k}.
\]
Orthogonality of the traced-out vectors removes cross terms. Orthogonality of the retained vectors then makes the remaining expression spectral:
\[
\boxed{\rho_A=\sum_k\lambda_k\proj{a_k},\qquad
\rho_B=\sum_k\lambda_k\proj{b_k}.}
\]
Thus \(\rho_A\ket{a_k}=\lambda_k\ket{a_k}\), and similarly for \(B\). The nonzero reduced spectra agree, even when the subsystem dimensions differ. Schmidt amplitudes are \(\sqrt{\lambda_k}\); their squares are the marginal eigenvalues. Both local Schmidt measurements reveal the same perfectly distinguishable label \(k\), with probability \(\lambda_k\).

\lead{Counting arbitrary terms would give false entanglement}
The product state
\[
\ket0\ket+=\frac{\ket0\ket0+\ket0\ket1}{\sqrt2}
\]
has two equal terms but only one Schmidt term, \(\ket0\ket+\). Its two displayed \(A\)-vectors coincide. Treating \(1/2,1/2\) as reduced eigenvalues would incorrectly assign one ebit to a product state. A valid expansion therefore need not be Schmidt.

\lead{Construct Schmidt form by diagonalizing a marginal}
Diagonalize \(\rho_A\), keeping its positive eigenvalues and corresponding orthonormal eigenvectors. Define
\[
\ket{b_k}=\frac{(\bra{a_k}\otimes\id_B)\ket\Psi}{\sqrt{\lambda_k}},
\qquad
\braket{b_j|b_k}
=\frac{\bra{a_k}\rho_A\ket{a_j}}{\sqrt{\lambda_j\lambda_k}}
=\delta_{jk}.
\]
These partner vectors reconstruct \(\Psi\) and are automatically orthonormal. Every pure bipartite state therefore admits Schmidt form. Equivalently, it is the singular-value decomposition of the ket's coefficient matrix; squared singular values give the reduced eigenvalues.

\takeaway{For Shor, \(\lambda_c=M_c/Q\), \(\ket{a_c}=\ket{C_c}\), and \(\ket{b_c}=\ket{e_c}\). Rank \(r>1\) proves entanglement because the coherent input pairs orthogonal combs with distinct modular values. The input/output and eigenphase expressions describe this same state in different work bases.}

\phead{SCHMIDT BASES AND PHASE RAMPS}{5b. Which basis changes preserve Schmidt form?}{A fixed spectrum allows rotations only within equal-weight blocks.}
For distinct positive \(\lambda_k\), the paired Schmidt vectors are fixed up to reordering and compensating phases. Equal eigenvalues permit compensated unitary rotations within their common eigenspaces. Arbitrary changes of product basis still describe the state, but generally lose Schmidt form.

\lead{Equal comb weights give many Schmidt decompositions}
If \(r\mid Q\), all occupied eigenvalues are \(1/r\). For any unitary \(r\times r\) matrix \(V\), define
\[
\ket{A_s}=\sum_cV_{sc}\ket{C_c},
\qquad
\ket{B_s}=\sum_cV_{sc}^{*}\ket{e_c}.
\]
Both families are orthonormal, and unitarity gives
\[
\frac1{\sqrt r}\sum_s\ket{A_s}\ket{B_s}
=\frac1{\sqrt r}\sum_{c,d}
\underbrace{\left(\sum_sV_{sc}V_{sd}^{*}\right)}_{\delta_{cd}}
\ket{C_c}\ket{e_d}
=\ket{\Psi_2}.
\]
Choose \(V_{sc}=\omega^{sc}/\sqrt r\). The two rotated families are the counting ramps and work eigenstates:
\[
\ket{\Phi_s}=\frac1{\sqrt r}\sum_c\omega^{sc}\ket{C_c},
\qquad
\ket{u_s}=\frac1{\sqrt r}\sum_c\omega^{-sc}\ket{e_c}.
\]
Consequently, both the comb expansion and \(r^{-1/2}\sum_s\ket{\Phi_s}\ket{u_s}\) are Schmidt decompositions in this case. The opposite phase signs make the two rotations compensate.

\lead{Finite-length ramps need not be orthogonal}
For general \(Q\), the definition in Section 4 gives
\[
\boxed{\braket{\Phi_s|\Phi_{s'}}
=\frac1Q\sum_{x=0}^{Q-1}e^{2\pi\ii x(s'-s)/r}.}
\]
Complete periods cancel for \(s\ne s'\) when \(r\mid Q\). If \(r\nmid Q\), an incomplete final period remains. In particular, for \(r>1\),
\[
\braket{\Phi_s|\Phi_{s+1}}
=\frac{1-e^{2\pi\ii Q/r}}{Q(1-e^{2\pi\ii/r})}\ne0.
\]
Thus the family is not mutually orthogonal; some nonadjacent pairs may still be orthogonal. Orthogonal computational coordinates do not make every pair of superpositions orthogonal.

For a minimal illustration, let \(Q=4,r=3\), with \(\omega=e^{2\pi\ii/3}\). Then
\[
\ket{\Phi_0}=\tfrac12(1,1,1,1)^{\mathsf T},\qquad
\ket{\Phi_1}=\tfrac12(1,\omega,\omega^2,1)^{\mathsf T},
\qquad
\braket{\Phi_0|\Phi_1}=\tfrac14.
\]
Three contributions cancel; the fourth survives. Its reduced eigenvalues are \((1/2,1/4,1/4)\), rather than the equal ramp-ensemble weights \(1/3\).

\takeaway{The comb Schmidt form is exact for every \(Q\ge r\). If \(r\nmid Q\), rotations remain possible among combs with equal \(M_c\), but mixing different-weight blocks generally destroys Schmidt form. The phase expansion remains exact and useful for readout even when it is not spectral.}
